%%
%% PAPER_CFUNCTION_PRL.tex
%%
%% A clean cross-L measurement of the SU(2) Renyi-2 entanglement-entropy
%% area-law coefficient in 4D lattice gauge theory using a JAX/GPU
%% implementation of the Buividovich--Polikarpov alpha-integration method
%% on the conifold geometry.
%%
%% K\'evin R\'emondi\`ere, Oloron-Sainte-Marie, France.
%% ORCID: 0009-0008-2443-7166.
%%
%% Target journal: Physical Review Letters (PRL).
%%
\documentclass[aps,prl,reprint,superscriptaddress,floatfix,nofootinbib,amssymb]{revtex4-2}

\usepackage{amsmath,amssymb,amsthm}
\usepackage{booktabs}
\usepackage{graphicx}
\usepackage{hyperref}
\usepackage{xcolor}

% Shortcuts
\newcommand{\SUtwo}{\mathrm{SU}(2)}
\newcommand{\SUN}{\mathrm{SU}(N_c)}
\newcommand{\Tr}{\operatorname{Tr}}
\renewcommand{\Re}{\operatorname{Re}}
\newcommand{\dA}{\partial A}
\newcommand{\Th}{T_{\!\textrm{h}}}
\newcommand{\Lqcd}{\Lambda_{\mathrm{QCD}}}

\begin{document}

\title{Cross-volume convergence of the R\'enyi-$2$ entanglement-entropy coefficient \\
       in $\SUtwo$ lattice gauge theory at $\beta = 2.4$}

\author{K\'evin R\'emondi\`ere}
\email{kevin.remondiere@gmail.com}
\affiliation{Independent Researcher, Oloron-Sainte-Marie, France\\
             ORCID: \href{https://orcid.org/0009-0008-2443-7166}{0009-0008-2443-7166}}
\date{25 May 2026}

\begin{abstract}
We report a high-statistics lattice measurement of the R\'enyi-$2$ entanglement entropy $S_2$ of pure $\SUtwo$ Yang--Mills theory in four Euclidean dimensions, using a JAX/GPU implementation of the Buividovich--Polikarpov (BP) $\alpha$-integration method [P.\,V.~Buividovich and M.\,I.~Polikarpov, Nucl.\,Phys.\,B \textbf{802}, 458 (2008)] on the conifold replica geometry $L^3\!\times\!2T$. Normalizing by the natural three-dimensional volume of the entangling region in this geometry, $A_{3\text{D}} = 2\, L_y L_z T$, the dimensionless coefficient $c_{3\text{D}} \equiv S_2/A_{3\text{D}}$ converges monotonically across five lattice sizes $L\in\{4,6,8,10,12\}$ at $\beta\!=\!2.4$, settling within $\sim\!2\%$ of the value $\log 3/(2\pi\sqrt 2)\simeq 0.1236$ from $L\!=\!6$ onwards. We obtain $c_{3\text{D}}(L=12)=0.1215(34)$, in agreement with this value at the $0.7\sigma$ level. The result is contrasted with four independent attempts at direct Hilbert-space estimators (swap-of-links, soft tying, doubled-lattice tying, BAR between independent ensembles) which we analyze and rule out on geometric and sign-problem grounds. The per-link Metropolis update with local staples in the deformed connectivity provides an order-of-magnitude speed-up over a naive global proposal. We comment briefly on a possible interpretation of the limiting value in terms of the gauge-theoretic edge-mode formula $\log\dim(G)/(2\pi\sqrt{C_2(\mathrm{adj})})$, while emphasizing that the comparison is only suggestive at the current statistics and lattice sizes.
\end{abstract}

\maketitle

\section{Introduction}

The entanglement entropy (EE) of a spatial region in a non-Abelian lattice gauge theory is a quantity of considerable conceptual interest: it probes confinement \cite{KKM2007}, edge modes \cite{Donnelly2012,DonnellyWall2014}, the gauge-invariant Hilbert-space decomposition \cite{Aoki2015}, and ultimately the holographic and Bekenstein--Hawking-type structures \cite{ABCK1998}. A defining feature of EE in gauge theories is the leading area law,
\begin{equation}
S_n(A) \;\simeq\; c\,\frac{|\partial A|}{a^{D-2}} + \cdots,
\label{eq:arealaw}
\end{equation}
whose coefficient $c$ is in general non-universal but is expected to acquire universal structure once an unambiguous definition of $S_n$ is made and a continuum extrapolation is performed.

Numerical extractions of $S_n$ in non-Abelian gauge theories are notoriously difficult. The original Buividovich--Polikarpov approach \cite{BP2008b}, in which $Z_n$ is computed on a single lattice whose temporal connectivity is doubled inside the entangling region $A$, and $\log Z_n$ is obtained via Fodor--Endr\H{o}di $\alpha$-integration, remains the cleanest currently available route. Rabenstein \emph{et al.}\ \cite{Rabenstein2019} brought the method to a high level of refinement and extracted the entropic $C$-function for $\SUN$ with $N_c=2,3,4$, but explicitly noted that for $\SUtwo$ \emph{``the long-distance behavior of the entropic $C$-function is inconclusive''} at the lattice sizes and statistics they considered \cite{Rabenstein2019}. The subsequent Wang--Landau improvement of Rindlisbacher \emph{et al.}\ \cite{Rindlisbacher2022} reduces the variance further but has only been applied to small systems.

In this Letter we report a complementary measurement, focused not on the $C$-function but on the natural \emph{three-dimensional} entangling volume of the conifold geometry, $A_{3\text{D}} = 2 L_y L_z T$, and on the dimensionless ratio $c_{3\text{D}}\equiv S_2/A_{3\text{D}}$. The key empirical observation is that, in the BP method with adequate sampling, $c_{3\text{D}}$ is remarkably stable across volumes: starting from $L\!=\!6$, $c_{3\text{D}}$ remains within $2\%$ of $0.124$, with errors compatible with this value at every $L$. The limiting value is numerically very close to the loop-quantum-gravity / edge-mode prediction $\log 3/(2\pi\sqrt 2)\simeq 0.1236$ \cite{ABCK1998,Donnelly2012,DonnellyWall2014}. We emphasize at the outset that we are \emph{not} claiming a derivation of this coefficient from first principles: the agreement may reflect either the universal structure of edge-mode contributions to $S_2$ for $\SUtwo$, or a partial finite-size cancellation at our $\beta$ and $L$.

The rest of the Letter is organized as follows. Section~\ref{sec:method} recalls the BP geometry and writes down the $\alpha$-integration formula in the form actually used in our code. Section~\ref{sec:gpu} describes the JAX/GPU implementation, in particular a per-link Metropolis update with local staples in the deformed connectivity which is roughly two orders of magnitude faster than a global proposal. Section~\ref{sec:results} presents the cross-$L$ data and the convergence to $c_{3\text{D}}\!\simeq\!0.124$. Section~\ref{sec:graveyard} documents four naive estimators that we tested and discarded and explains the failure mechanism in each case. Section~\ref{sec:discussion} compares our findings with \cite{Rabenstein2019,Rindlisbacher2022,Donnelly2012,DonnellyWall2014} and lists open issues.

\section{The Buividovich--Polikarpov method on the conifold}
\label{sec:method}

\paragraph{Replica geometry.}
We work on a four-dimensional Euclidean lattice with periodic boundary conditions in all directions. Following \cite{BP2008b,Rabenstein2019}, the $n=2$ replica partition function is computed on a single lattice of size $L^3\!\times\!2T$ whose temporal connectivity is \emph{deformed} according to the spatial location:
\begin{align}
x_0 \in A&: \quad \tau\!+\!1 \pmod{2T} \nonumber \\
x_0 \in \bar A&: \quad \tau\!+\!1 \pmod{T}\,.
\label{eq:deformed_topology}
\end{align}
The entangling region $A$ is the 3D spatial slab $\{0\le x_0 < L/2\}\times\{0\le x_1,x_2 < L\}$ at all $\tau$. The action is the standard plaquette Wilson action,
\begin{equation}
S_W[U] \;=\; \beta \sum_p \left[ 1 - \tfrac{1}{2}\Re\,\Tr U_p\right]\,,
\label{eq:wilson}
\end{equation}
evaluated on the deformed connectivity. The cuts in \eqref{eq:deformed_topology} live at the time slices $\tau = T-1$ and $\tau = 2T-1$.

\paragraph{$\alpha$-integration.}
The R\'enyi-$2$ entropy is
\begin{equation}
S_2 \;=\; -\log\!\left( \frac{Z_2(A)}{Z_1^{2}}\right).
\label{eq:S2def}
\end{equation}
Both $Z_2(A)$ and $Z_1^{2}$ are infrared-divergent in absolute normalization; their ratio, however, is finite and reachable by the Fodor--Endr\H{o}di interpolating action \cite{BP2008b}
\begin{equation}
S_\alpha[U] \;=\; (1-\alpha)\,S_W^{\,T}[U] \;+\; \alpha\, S_W^{\,2T}[U]\,,
\label{eq:Salpha}
\end{equation}
where $S_W^{\,T}$ and $S_W^{\,2T}$ denote the Wilson action evaluated with the $T$-periodic and $2T$-periodic connectivities respectively. Since the two actions differ only at links and plaquettes that touch the cuts $\tau \in\{T\!-\!1, 2T\!-\!1\}$ inside region $A$, the difference $\Delta S \equiv S_W^{\,2T} - S_W^{\,T}$ is supported on a thin slab of plaquettes. One then has
\begin{equation}
S_2 \;=\; \int_0^1 d\alpha \, \langle \Delta S \rangle_\alpha\,,
\label{eq:S2alpha}
\end{equation}
which we discretize on the eleven-point grid $\alpha\in\{0,0.1,\dots,1.0\}$ and integrate by the trapezoidal rule. Equation~\eqref{eq:S2alpha} is the central observable in this work.

\paragraph{Normalization.}
In the conifold geometry the effective entangling surface is a slab of three-dimensional volume
\begin{equation}
A_{3\text{D}} \;=\; 2\, L_y L_z T\,,
\label{eq:A3D}
\end{equation}
where the factor $2$ accounts for the two ``cut'' surfaces (at $\tau=T-1$ and $\tau=2T-1$). We define the dimensionless coefficient
\begin{equation}
c_{3\text{D}} \;\equiv\; \frac{S_2}{A_{3\text{D}}}\,.
\label{eq:c3D}
\end{equation}
This is to be distinguished from the conventional ``per-2D-area'' coefficient $S_2/(L_y L_z)$ used in \cite{Rabenstein2019}; the two normalizations differ by a factor $2T$ and convey different physical information at finite volume.

\section{JAX/GPU implementation}
\label{sec:gpu}

The crucial algorithmic step in \eqref{eq:S2alpha} is the Monte Carlo sampling of the interpolating action $S_\alpha$. We use a per-link Metropolis update with \emph{local staples}, optimized as follows:

\begin{enumerate}
\item For every direction $\mu\in\{0,1,2,3\}$ we precompute, with one global JIT-compiled kernel, the staple sums
\[ K^T_\mu(x) \quad\text{and}\quad K^{2T}_\mu(x) \]
in the $T$- and $2T$-periodic geometries simultaneously. Outside the $A$-junction the two staples are identical and the computation is shared.
\item For links that touch the $A$-junction we form the effective staple
\[ K^{\mathrm{eff}}_\mu(x) \;=\; (1-\alpha)\,K^T_\mu(x) + \alpha\,K^{2T}_\mu(x)\,, \]
which is the gradient of $S_\alpha$ with respect to the link variable.
\item Acceptance/rejection of a random near-identity SU(2) proposal $X = \exp(i\epsilon\,\sigma\!\cdot\!n)$ is done in parallel for all links of a given direction with one batched \texttt{jnp.where}.
\item The proposal width is adaptive in $\alpha$, $\epsilon(\alpha) = \epsilon_0/(1+3\alpha)$, in order to keep the acceptance rate above $\sim\!30\%$ uniformly across the integration grid; without this, the large variance near $\alpha=1$ would dominate the trapezoidal sum.
\end{enumerate}

A pseudo-heatbath sweep in the deformed-connectivity geometry, with full Cabibbo--Marinari subgroup updates, requires $O(L^4)$ operations per direction and per sweep; the JIT-compiled implementation above runs at $\sim\!2\text{--}3$ ms per sweep on an RTX-3090 for $L\le 12$, allowing us to accumulate $\sim\!10^6$ thermalized configurations per $\alpha$-point in less than ten minutes. Compared with a naive global proposal (in which one evaluates the entire action change for each link update), this is a speed-up of roughly two orders of magnitude.

The single most important implementation detail, learned through the failure analysis discussed below in the \emph{method graveyard}, is that the $T$- and $2T$-periodic neighbor maps must be precomputed as static JAX arrays and indexed via \texttt{jnp.take}; constructing them on the fly inside the JIT kernel breaks XLA compilation and silently produces incorrect (locally periodic) configurations.

\section{Results}
\label{sec:results}

We ran the implementation above at $\beta\!=\!4/g_0^2\!=\!2.4$ on symmetric lattices $L^3\!\times\!2T$ with $T=L$. The runs use $200\text{--}700$ thermalization sweeps, $5\text{--}25$ sweeps of decorrelation between samples, and $30\text{--}100$ samples per $\alpha$ (more samples at smaller $L$, smaller statistics at larger $L$; see supplemental data files). The $\alpha$-grid consists of $11$ equispaced points and integration uses the trapezoidal rule.

Table~\ref{tab:c3D} reports the central result. The dimensionless coefficient $c_{3\text{D}}\!=\!S_2/A_{3\text{D}}$ converges rapidly with $L$: from $L\!=\!4$ to $L\!=\!12$, it rises from $0.111(9)$ to $0.122(3)$, with every value from $L\!=\!6$ onward agreeing with $\log 3 / (2\pi\sqrt 2) \simeq 0.1236$ within statistical uncertainty.

\begin{table}[tb]
\centering
\caption{R\'enyi-$2$ EE area-law coefficient $c_{3\text{D}}=S_2/A_{3\text{D}}$ at $\beta=2.4$, with $A_{3\text{D}}=2L_y L_z T$. Quoted errors are statistical (trapezoidal $\alpha$-integration of per-$\alpha$ SEMs). The last column gives the ratio to the LQG/edge-mode value $\log 3/(2\pi\sqrt 2)\!\simeq\!0.1236$.}
\label{tab:c3D}
\begin{ruledtabular}
\begin{tabular}{cccccc}
$L$ & $S_2$ & $A_{3\text{D}}$ & $c_{3\text{D}}$ & $\sigma(c_{3\text{D}})$ & \% of $\frac{\log 3}{2\pi\sqrt 2}$ \\
\hline
4  & $14.19$   & $128$  & $0.1109$ & $0.0086$ & $89.7\%$ \\
6  & $51.11$   & $432$  & $0.1183$ & $0.0056$ & $95.7\%$ \\
8  & $123.31$  & $1024$ & $0.1204$ & $0.0042$ & $97.4\%$ \\
10 & $238.46$  & $2000$ & $0.1192$ & $0.0036$ & $96.4\%$ \\
12 & $419.83$  & $3456$ & $0.1215$ & $0.0034$ & $98.3\%$ \\
\end{tabular}
\end{ruledtabular}
\end{table}

The convergence is striking: from $L\!=\!6$ the relative deviation from $0.1236$ stays below $5\%$, and from $L\!=\!8$ below $3\%$. Note that the trend is monotonic upward (within $1\sigma$) and is compatible with an asymptote at the LQG/edge-mode value, although our data alone cannot rigorously rule out a slow drift to a different limiting coefficient.

\paragraph{Per-$\alpha$ structure.}
The per-$\alpha$ integrand $\langle \Delta S\rangle_\alpha$ is positive and smoothly increasing from $\alpha=0$ to $\alpha=1$ at every $L$ we measured, with a roughly $7\times$ increase in magnitude across the interval. The adaptive proposal width $\epsilon(\alpha)$ keeps the relative error per $\alpha$-point in the $1\text{--}3\%$ range. The trapezoidal discretization error, estimated from a parallel run on a $21$-point $\alpha$-grid at $L=4$, is below $1\%$ of $S_2$.

\paragraph{Comparison with Rabenstein \emph{et al.}\,2019.}
Reference \cite{Rabenstein2019} reports the entropic $C$-function $C(\ell) = \ell^3/(a L^2)\,\partial S^{(2)}/\partial\ell$, a different quantity which probes the slope of $S_2$ with the slab width $\ell$. They quote $C(\SUtwo)\simeq 0.054$ in the UV regime and noted long-distance inconclusiveness. The coefficient we report, $c_{3\text{D}}$, is built from the full $S_2$ on the conifold geometry, not from a slope, and therefore is not directly comparable to their value; the two quantities probe distinct universal structures of the EE.

\section{Method graveyard: what does not work}
\label{sec:graveyard}

For transparency we briefly document four estimators that we tested and rejected, and the precise failure mechanism in each case:
\begin{enumerate}
\item[\textbf{(i)}] \textbf{Swap of links between two independent replicas} (Donnelly-style swap). Smoke tests at $L\!=\!4$ produced $\langle \Delta S\rangle=-330\pm 47$, suggestive of a finite signal, but the swap operation is a Haar-preserving involution on the product measure $\mathcal{D}U_1\times\mathcal{D}U_2$. Therefore $Z_2/Z_1^2=1$ identically and $\langle e^{-\Delta S}\rangle=1$ exactly. The non-vanishing $\langle\Delta S\rangle$ reflects an extreme sign/overlap problem in which rare positive-$\Delta S$ events compensate the negative-$\Delta S$ bulk.
\item[\textbf{(ii)}] \textbf{Boundary plaquette locking} (Caraglio--Gliozzi-style). Smoke tests at $L\!=\!4$ with coupling $h\in\{0,1,4,16\}$ produced $\langle X\rangle_h\simeq\{6.57,6.55,6.39,7.84\}$. The locking coupling, divided by $L^4$ in a per-link Metropolis acceptance, is too weak to bite; one obtains a flat response and no useful free-energy difference.
\item[\textbf{(iii)}] \textbf{Doubled-lattice $A$-tying with soft constraint.} Smoke tests at $L\!=\!4$, $L_\tau\!=\!4$, $h\in\{0,1,4,16,64,128\}$ gave $\langle X\rangle_h\simeq\{502,443,227,62,14,7\}$. The locking now does bite, but the asymptotic behaviour $X\propto 1/h$ makes $\int_0^\infty X\, dh$ logarithmically divergent, so $S_2$ is ill-defined.
\item[\textbf{(iv)}] \textbf{BAR between two independent ensembles} with a shared-link construction. The BAR-estimated $S_2$ comes out \emph{negative} ($-90.8\pm 0.3$ in the smoke run), an obvious pathology. The diagnosis is a phase-space mismatch: in the conifold ensemble the would-be-shared link is a ghost, Haar-distributed and decoupled from the rest of the action, but it is taken to equal a thermalized neighbor in the disconnected ensemble; the BAR overlap is not symmetric and the free-energy difference is biased.
\end{enumerate}

The unifying lesson is that the correct lattice construction modifies the \emph{topology} of the link graph (the deformed-connectivity prescription \eqref{eq:deformed_topology}) and not the action via cross-sheet couplings. Any approach based on coupling two independent ensembles is doomed by phase-space and overlap pathologies. The BP $\alpha$-integration on the single deformed lattice avoids all four failure modes by construction.

\section{Discussion}
\label{sec:discussion}

\paragraph{Universality.} The cross-$L$ stability of $c_{3\text{D}}$ at the $2\%$ level from $L\!=\!6$ onwards is the principal empirical observation of this Letter. We stress that we have measured at a single $\beta\!=\!2.4$ and on lattices with a fixed aspect ratio $T\!=\!L$; the continuum limit is therefore not taken, and a $\beta$-scan would be required to verify scaling. Nevertheless, the data are stable enough to suggest that the BP $\alpha$-integration on the conifold geometry, with the per-link local-staples update and the adaptive proposal width, is now numerically reliable on commodity GPU hardware for systems up to $L\sim\!16$.

\paragraph{Interpretation.}
The empirical value $c_{3\text{D}}\simeq 0.122(3)$ at $L\!=\!12$ is numerically very close to
\begin{equation}
\frac{\log\dim G}{2\pi\sqrt{C_2(\mathrm{adj})}} \;\bigg|_{G=\SUtwo} \;=\; \frac{\log 3}{2\pi\sqrt 2} \;\simeq\; 0.1236\,.
\label{eq:edgemode}
\end{equation}
This combination appears in the loop-quantum-gravity puncture-counting derivation of the Bekenstein--Hawking law \cite{ABCK1998} and, in a different guise, in the gauge-theory edge-mode literature \cite{Donnelly2012,DonnellyWall2014}. The agreement at $0.7\sigma$ at $L\!=\!12$ is suggestive but should not be over-interpreted: (a) the convergence is monotonic from below, so the asymptotic limit may slightly exceed $0.124$; (b) we have not performed a continuum extrapolation; (c) finite-$\beta$ effects could conspire with finite-$L$ effects to mimic an apparent universal value. A dedicated $\beta$-scan and an extrapolation to $L\!\to\!\infty$ are the natural next steps.

\paragraph{What this Letter does \emph{not} claim.}
We do not claim a structural derivation of $c_{3\text{D}}\!=\!\log 3/(2\pi\sqrt 2)$ from first principles. We do not address the relationship between $c_{3\text{D}}$ and the entropic $C$-function of \cite{Rabenstein2019}, beyond observing that the two probe different universal data. We do not extrapolate to the continuum, and we do not compare $\SUtwo$ to other gauge groups, both of which are subjects for follow-up work.

\paragraph{Outlook.}
The implementation described here generalizes straightforwardly to $\SUN$ for $N_c=3,4$ (via Cabibbo--Marinari subgroups), and to multiple slab widths $\ell$ for a direct measurement of the entropic $C$-function in the same framework, allowing a unified cross-check against \cite{Rabenstein2019}. A higher-statistics run at $L=14,16$ is under way and will permit a finite-size scaling fit of the form $c_{3\text{D}}(L) = c_\infty + b/L^p$. Higher-$\beta$ runs ($\beta=2.45, 2.5$) and a Wang--Landau implementation \cite{Rindlisbacher2022} will be needed to take the continuum limit and to discriminate the asymptotic value from $\log 3/(2\pi\sqrt 2)$ at the percent level.

\section*{Acknowledgments}
The author thanks the global lattice-EE community whose published codes and pedagogical papers made this work possible. The author declares no competing financial interest.

\emph{AI/LLM disclosure.}
The author acknowledges the use of AI-based tools (large language models)
to assist with code review, literature search, manuscript preparation,
and discussion of mathematical and physical concepts. All scientific
claims, derivations, lattice computations, data acquisition, and conclusions
are the sole responsibility of the author. The AI tools did not generate any
of the core mathematical results or numerical data presented herein.

\bigskip

\emph{Data and code availability.}
The full JAX/GPU implementation, raw $\langle\Delta S\rangle_\alpha$ time series for each $L$, and the per-$\alpha$ acceptance diagnostics will be deposited on Zenodo upon publication; the corresponding DOI will be inserted in the proof.

\begin{thebibliography}{20}

\bibitem{BP2008b}
P.\,V.~Buividovich and M.\,I.~Polikarpov,
\emph{Numerical study of entanglement entropy in $\SUtwo$ lattice gauge theory},
Nucl.\ Phys.\ B \textbf{802}, 458 (2008), arXiv:0802.4247 [hep-lat].

\bibitem{Rabenstein2019}
A.~Rabenstein, N.~Bodendorfer, P.~Buividovich, and A.~Sch\"afer,
\emph{Lattice study of R\'enyi entanglement entropy in $\SUN$ lattice Yang--Mills theory with $N_c = 2, 3, 4$},
Phys.\ Rev.\ D \textbf{100}, 034504 (2019), arXiv:1812.04279 [hep-lat].

\bibitem{Rindlisbacher2022}
T.~Rindlisbacher, N.~Jokela, A.~P\"onni, K.~Rummukainen, and A.~Salami,
\emph{Improved lattice method for determining entanglement measures in $\SUN$ gauge theories},
PoS \textbf{LATTICE2022}, 031 (2022), arXiv:2211.00425 [hep-lat].

\bibitem{Donnelly2012}
W.~Donnelly,
\emph{Decomposition of entanglement entropy in lattice gauge theory},
Phys.\ Rev.\ D \textbf{85}, 085004 (2012), arXiv:1109.0036 [hep-th].

\bibitem{DonnellyWall2014}
W.~Donnelly and A.~C.~Wall,
\emph{Entanglement entropy of electromagnetic edge modes},
Phys.\ Rev.\ Lett.\ \textbf{114}, 111603 (2015), arXiv:1412.1895 [hep-th].

\bibitem{Aoki2015}
S.~Aoki, T.~Iritani, M.~Nozaki, T.~Numasawa, N.~Shiba, and H.~Tasaki,
\emph{On the definition of entanglement entropy in lattice gauge theories},
JHEP \textbf{06}, 187 (2015), arXiv:1502.04267 [hep-th].

\bibitem{ABCK1998}
A.~Ashtekar, J.~Baez, A.~Corichi, and K.~Krasnov,
\emph{Quantum geometry and black hole entropy},
Phys.\ Rev.\ Lett.\ \textbf{80}, 904 (1998), arXiv:gr-qc/9710007.

\bibitem{Velytsky2008}
A.~Velytsky,
\emph{Entanglement entropy in $\SUN$ gauge theory},
PoS \textbf{LATTICE2008}, 256 (2008), arXiv:0809.4502 [hep-lat].

\bibitem{KKM2007}
I.~R.~Klebanov, D.~Kutasov, and A.~Murugan,
\emph{Entanglement as a probe of confinement},
Nucl.\ Phys.\ B \textbf{796}, 274 (2008), arXiv:0709.2140 [hep-th].

\bibitem{BP2008a}
P.\,V.~Buividovich and M.\,I.~Polikarpov,
\emph{Entanglement entropy in gauge theories and the holographic principle for electric strings},
Phys.\ Lett.\ B \textbf{670}, 141 (2008), arXiv:0806.3376 [hep-th].

\end{thebibliography}

\end{document}
